\section{The actual kernel layer, its derivative image, and its quotient}
\label{sec:kernel-layer-continuation}

This section integrates the complete kernel-layer delivery of
12 September 2026.  Its archive hash and member hashes are recorded in
\path{sources/web_kernel_layer_delivery/LOCAL_STAGING_PROVENANCE.json}.
The source note's formulas (9)--(28) identify the interpolation kernel
with an explicit basis of the next original theta relation layer.
We first prove those maps in the retained coordinates and then calculate
the exact derivative image and the entire quotient of that layer.
The new quotient calculation retains the part of the layer not reached
by the summed Euler derivative.

\subsection{The original arithmetic source and its constrained minimum}

Keep \(g=2\xi\), a nonempty full packet
\(h(s)=\prod_{\rho\in Z}(s-\rho)^{m_\rho}\), and
\(d=\deg h\).  Fix \(k\geq1\), and set
\[
 E=E_h^{\otimes k},\quad E_h=\mathbb C[s]/(h),\quad
 A=\sum_{j=1}^k A_{h,j},\quad
 v_h=g/h,\quad \upsilon_h=j_hv_h\in E_h^\times.
                                                               \tag{KL.1}
\]
All tensor coordinates use the original monomial remainders.  The
function \(F_h\), with Mellin transform \(g/h\), is the original
test function satisfying \(h(D)F_h=f_0\).  The exact source map is
\[
 \mathcal T(P)=P(D_1,\ldots,D_k)F_h^{\otimes k},\quad
 \mathcal M_k\mathcal T(P)
       =P(s_1,\ldots,s_k)\prod_{j=1}^k v_h(s_j),
 \qquad s_j=\tfrac12+it_j.                                \tag{KL.2}
\]
The Mellin multiplier proves injectivity on polynomials.  Its Hilbert
norm is the polynomial norm for the unchanged product measure
\[
 d\nu_h(t)=|v_h(\tfrac12+it)|^2\frac{dt}{2\pi},\qquad
 \nu_h^{\otimes k}.                                      \tag{KL.3}
\]
Let \(\mathcal P_N\) be the polynomials of total degree at most
\(N\), and let \(\mathcal J_N:\mathcal P_N\to E\) be full
reduction by \(h(s_1),\ldots,h(s_k)\), followed by multiplication
by the original unit \(\upsilon_h^{\otimes k}\).
For \(N\geq k(d-1)\) this is onto, since the monomial remainder
basis has at most that total degree.  Fix an integer
\(N\geq k(d-1)\) for the remainder of this section; all subsequent
levels \(N+1,N+2\) satisfy the same threshold.
Its kernel and its original
theta image are
\[
 \mathcal I_N=\mathcal P_N\cap(h(s_1),\ldots,h(s_k)),
 \qquad \mathcal L_N=\mathcal T\mathcal I_N.              \tag{KL.4}
\]

Let \(p_j(s)\) be the actual monic orthogonal polynomial of
\(\nu_h\), and retain \(\kappa_j=\|p_j\|_{\nu_h}^2>0\).
All moments exist by the proved arithmetic decay.  The Gram is
strictly positive because the density is positive except at a
discrete set; this proves existence and uniqueness by finite
orthogonal projection.  For multi-indices set
\[
 p_\alpha=\prod_jp_{\alpha_j},\quad
 \kappa_\alpha=\prod_j\kappa_{\alpha_j},\quad
 a_\alpha=\bigotimes_j\bigl(\upsilon_h[p_{\alpha_j}]_h\bigr).
                                                               \tag{KL.5}
\]
The monic triangular change preserves the full space of degree at
most \(N\), and product integration makes its Gram diagonal.
It follows, by direct constrained minimization, that the original
canonical representative and its metric are
\[
 \mathcal K_N=\sum_{|\alpha|\leq N}
                 \frac{a_\alpha a_\alpha^*}{\kappa_\alpha}>0,
 \qquad G_N=\mathcal K_N^{-1},
 \qquad
 R_Nu=\sum_{|\alpha|\leq N}\mathcal T(p_\alpha)
                 \frac{a_\alpha^*G_Nu}{\kappa_\alpha}.
                                                               \tag{KL.6}
\]
Indeed the full jet of this vector is \(\mathcal K_NG_Nu=u\).
If a coefficient vector has zero full jet, its inner product with
the displayed coefficient vector is zero.  Every other representative
of that same jet is its sum with such a vector; its squared norm is
therefore \(u^*G_Nu\) plus the squared norm of the added relation.
This proves the unique minimum, its metric, and
\(R_N\perp\mathcal L_N\).  These are the original maps
\(J^{(k)}R_N=I\) and \(q^{(k)}R_N=\sigma\).
The notation \(\mathcal K_N\) here includes the local unit.
To specify the comparison in the full tensor dimension, put
\[
 r_\alpha=\bigotimes_j[p_{\alpha_j}]_h,\qquad
 K_N^{\rm rem}=\sum_{|\alpha|\leq N}
             \frac{r_\alpha r_\alpha^*}{\kappa_\alpha}.
\]
Then \(\mathcal K_N=U_{\upsilon_h^{\otimes k}}K_N^{\rm rem}
U_{\upsilon_h^{\otimes k}}^*\), as an equality of operators on
\(E_h^{\otimes k}\).  When \(k=1\), \(K_N^{\rm rem}\) is
exactly the earlier remainder-only kernel \(K_N\).

\subsection{The complete next-layer frame and its theta primitives}

Write \(\Sigma_n=\{\alpha\in\mathbb N^k:|\alpha|=n\}\)
and \(\mathsf H_n=\mathbb C^{\Sigma_n}\).  The coordinate
identification sends \(x\in\mathsf H_n\) to the homogeneous
polynomial \(\sum_\alpha x_\alpha z^\alpha\); these \(z_j\)
are coefficient variables and are not substituted for the original
Mellin variables \(s_j\).
For \(\beta\in\Sigma_{N+1}\), define the actual functions
\[
 e_\beta=\mathcal T(p_\beta)-R_Na_\beta,
 \qquad \mathscr E_N x=\sum_\beta x_\beta e_\beta,
 \qquad
 \mathcal E_N=\mathcal L_{N+1}\cap\mathcal L_N^\perp.
                                                               \tag{KL.7}
\]
The full jet of \(e_\beta\) is zero, hence it is in
\(\mathcal L_{N+1}\).  The new-degree polynomial is orthogonal
to all old polynomials and \(R_Na_\beta\perp\mathcal L_N\),
so \(e_\beta\in\mathcal E_N\).
The leading total-degree terms of the \(e_\beta\) are the distinct
monomials \(s^\beta\), proving independence by the injectivity of
\(\mathcal T\).  For any element of \(\mathcal E_N\), subtract
its full degree-\(N+1\) combination of the \(e_\beta\).  The
remainder has zero jet and degree at most \(N\), hence is in
\(\mathcal L_N\).  It is also orthogonal to that space, so it is
zero.  This proves the exact isomorphism
\[
 \mathscr E_N:\mathsf H_{N+1}\xrightarrow{\sim}\mathcal E_N
       \xrightarrow{\ e\mapsto[e]\ }
                            \mathcal L_{N+1}/\mathcal L_N.
                                                               \tag{KL.8}
\]
The inverse of the second arrow is \([f]\mapsto(I-P_{\mathcal L_N})f\).

Every vector in (KL.7) has a primitive in the original tensor complex.
Its numerator polynomial is
\[
 p_\beta-\sum_{|\alpha|\leq N}
           p_\alpha\frac{a_\alpha^*G_Na_\beta}{\kappa_\alpha}.
                                                               \tag{KL.9}
\]
Successive monic division in the ordered variables
\(s_1,\ldots,s_k\) expresses this zero-jet polynomial as
\(\sum_i h(s_i)Q_{i,\beta}\), where
\(\deg Q_{i,\beta}\leq N+1-d\).  At each replacement the
leading power of the divided variable decreases and total degree
does not increase, proving the bound.  After the last division,
the remainder has degree below \(d\) in every variable; its full
jet is zero, so uniqueness of the complete remainder forces it to
vanish.  In the \(i\)-th tensor factor replace \(F_h\) by
\((-1)^{i-1}\phi_*\), and apply \(Q_{i,\beta}(D_1,\ldots,D_k)\)
to the whole tensor.  The preceding \(i-1\) degree-one factors
give the differential sign \((-1)^{i-1}\).  Together with
\(h(D)F_h=\Theta\phi_*\), the boundary of this primitive is
exactly \(+\mathcal T(h(s_i)Q_{i,\beta})\).  Summing proves
the asserted primitive with all ordered signs retained.

Put
\[
 \mathsf U_N=(a_\beta)_{\beta\in\Sigma_{N+1}},\qquad
 \mathsf D_{N+1}=\operatorname{diag}(\kappa_\beta),\qquad
 \mathsf H_N^{\rm lay}=\mathscr E_N^*\mathscr E_N.
\]
Expanding (KL.7) and using old/new polynomial orthogonality proves
\[
 \mathsf H_N^{\rm lay}
      =\mathsf D_{N+1}+\mathsf U_N^*G_N\mathsf U_N>0,
 \qquad \mathscr E_N^*R_N=-\mathsf U_N^*G_N.              \tag{KL.10}
\]
Define the projection loss \(Y_N=P_{\mathcal E_N}R_N\).
In particular its exact formula is
\(Y_N=-\mathscr E_N(\mathsf H_N^{\rm lay})^{-1}\mathsf U_N^*G_N\).
Its rank is exactly \(\operatorname{rank}\mathsf U_N\), because
every other displayed factor is injective or invertible on its
declared space.

\subsection{The exact incidence map and the derivative-image rank}

Define the actual old top-degree columns and their original norms by
\[
 \mathsf A_N=(a_\alpha)_{\alpha\in\Sigma_N}:\mathsf H_N\to E,
 \qquad \mathsf D_N=\operatorname{diag}(\kappa_\alpha)_{\Sigma_N}.
                                                               \tag{KL.11}
\]
Define the coordinate incidence map by
\[
 \mathsf T_n:\mathsf H_n\to\mathsf H_{n+1},\qquad
 (\mathsf T_nx)_\beta=
             \sum_{i:\beta_i>0}x_{\beta-\mathbf e_i}.
                                                               \tag{KL.12}
\]
On the specified homogeneous coefficient polynomials this is
multiplication by \(\ell(z)=z_1+\cdots+z_k\).  A product of two
nonzero polynomials over \(\mathbb C\) is nonzero: choose a
monomial order compatible with products and multiply the two leading
terms.  Since \(\ell\ne0\), \(\mathsf T_n\) is injective.

The full unprojected derivative defect and its projected layer map are
\[
 B_N=D^{(k)}R_N-R_NA,\qquad
 C_N=(I-P_{\mathcal L_N})B_N:E\to\mathcal E_N.
\]
They have the exact coordinate factorization
\[
 \boxed{C_N=\mathscr E_N\mathsf T_N\mathsf D_N^{-1}
                                  \mathsf A_N^*G_N.}    \tag{KL.13}
\]
To prove it, multiply the finite numerator in (KL.6) by
\(s_1+\cdots+s_k\).  Monicity makes its coefficient at the
new top polynomial \(p_\beta\) equal to
\(\sum_{i:\beta_i>0}a_{\beta-\mathbf e_i}^*G_Nu/
\kappa_{\beta-\mathbf e_i}\), precisely the row on the right
side of (KL.13).  The numerator of \(R_NAu\) has degree at
most \(N\).  Subtract the right side of (KL.13) from \(B_Nu\):
its numerator has degree at most \(N\) and its full jet is zero,
so it is an original relation in \(\mathcal L_N\).
Orthogonal projection proves (KL.13), including the original
quotient statement \([B_Nu]=[C_Nu]\) in
\(\mathscr B^{(k)}/\mathcal L_N\).

\begin{theorem}[Exact derivative rank and its full kernel]
With every original jet column retained,
\[
 \boxed{\operatorname{rank}C_N=\operatorname{rank}\mathsf A_N,
 \qquad \ker C_N=\mathcal K_N\ker\mathsf A_N^*.}         \tag{KL.14}
\]
Define the original weight form
\(W_N=A^*G_N+G_NA-kG_N\).  It obeys the refined bound
\[
 \operatorname{rank}W_N\leq
  2\min\{\operatorname{rank}\mathsf A_N,
                         \operatorname{rank}\mathsf U_N\}.
                                                               \tag{KL.15}
\]
\end{theorem}
\begin{proof}
Both \(\mathscr E_N\) and \(\mathsf T_N\) are injective;
\(\mathsf D_N\) and \(G_N\) are invertible.  Thus (KL.13)
has rank \(\operatorname{rank}\mathsf A_N^*\), which equals
the rank of its conjugate transpose.  Its kernel is
\(G_N^{-1}\ker\mathsf A_N^*=\mathcal K_N\ker\mathsf A_N^*\),
with the displayed invertible coordinate map retained.  The original
Green identity and orthogonality give
\(W_N=-(Y_N^*C_N+C_N^*Y_N)\).  Each summand has rank at most
the lesser of the two map ranks; the rank of a sum is at most the
sum of its ranks.  Use (KL.10) and (KL.14) to obtain (KL.15).
No nonzero homogeneous-leading-term assertion has been substituted
for an actual column \(a_\alpha\).
\end{proof}

\subsection{The whole quotient of the next layer by the derivative image}

For \(k\geq2\), let \(\mathsf H_n'\) be the homogeneous
polynomials of degree \(n\) in \(z_1,\ldots,z_{k-1}\).
There is a split exact sequence of vector spaces
\[
 0\longrightarrow\mathsf H_n
  \xrightarrow{\ \mathsf T_n\ }\mathsf H_{n+1}
  \xrightarrow{\ \rho_n\ }\mathsf H_{n+1}'
  \longrightarrow0,\qquad
 \rho_n f=f(z_1,\ldots,z_{k-1},-\!\sum_{i<k}z_i).
                                                               \tag{KL.16}
\]
For a complete proof and inverse, divide \(f\) by the monic
linear polynomial \(\ell=z_k+\sum_{i<k}z_i\) in the last
variable.  Every replacement retains total degree, giving uniquely
\[
 f=\ell q+\iota r,\qquad q\in\mathsf H_n,
 \quad r=\rho_nf\in\mathsf H_{n+1}',                    \tag{KL.17}
\]
where \(\iota\) includes a polynomial independent of \(z_k\).
For uniqueness set \(z_k=-\sum_{i<k}z_i\) first to obtain
\(r=0\), then use injectivity of multiplication by \(\ell\).
This also proves surjectivity, exactness and the specified section
\(\iota\).  The ordered last coordinate and the signs in the
substitution remain part of the map.

\begin{theorem}[Exact unused-layer quotient]
Set \(Z_N=\operatorname{im}(\mathsf D_N^{-1}\mathsf A_N^*)
\subseteq\mathsf H_N\).  Then the actual quotient is
\[
 \boxed{\mathcal E_N/\operatorname{im}C_N
   \simeq(\mathsf H_N/Z_N)\oplus\mathsf H_{N+1}'.}        \tag{KL.18}
\]
For an original layer representative \(e=\mathscr E_N f\),
divide \(f=\ell q+\iota r\) by (KL.17); the map sends
\([e]\) to \(([q],r)\).  Its inverse sends
\(([q],r)\) to \([\mathscr E_N(\ell q+\iota r)]\).
\end{theorem}
\begin{proof}
The invertibility of \(G_N\) in (KL.13) gives
\(\operatorname{im}C_N=\mathscr E_N\mathsf T_NZ_N\).
Changing \(e\) by such an element changes exactly \(q\) by an
element of \(Z_N\), while leaving \(r\) unchanged.  Thus the
first map is well-defined.  Changing the representative of \([q]\)
has the same effect on the inverse, so the inverse is well-defined.
The unique division (KL.17) proves both inverse identities.
For \(k=1\), \(\mathsf T_N\) is an isomorphism between the
one-dimensional shell spaces and the same statement holds with
\(\mathsf H_{N+1}'=0\).
\end{proof}

In particular \(\dim\mathsf H_{N+1}'=\binom{N+k-1}{k-2}\)
for \(k\geq2\), by counting monomials; the first summand has
dimension \(\binom{N+k-1}{k-1}-\operatorname{rank}\mathsf A_N\).
Their sum is exactly the original layer dimension minus
\(\operatorname{rank}C_N\).  The quotient therefore retains both
the full-jet rank loss and the explicitly complementary relation
directions.  For every coordinate vector, (KL.9) supplies its original
theta primitive.  The quotient is of actual test-function relations,
with those primitives, rather than only a quotient of dimensions.

\subsection{The derivative between consecutive relation quotients}

The original Euler sum has
\(D^{(k)}\mathcal L_N\subseteq\mathcal L_{N+1}\), because
it multiplies the numerator polynomial by \(s_1+\cdots+s_k\)
and preserves the ideal in (KL.4).  Hence it induces the well-typed
map
\[
 \partial_N:\mathcal L_{N+1}/\mathcal L_N
       \longrightarrow\mathcal L_{N+2}/\mathcal L_{N+1},
 \qquad [f]\longmapsto[D^{(k)}f].                       \tag{KL.19}
\]
Under the exact layer coordinates of (KL.8), its matrix is
\[
 \boxed{\partial_N=\mathscr E_{N+1}\mathsf T_{N+1}
                                           \mathscr E_N^{-1},}
                                                               \tag{KL.20}
\]
where each \(\mathscr E\) is followed or preceded by its quotient
isomorphism as in (KL.8).  Write \(e_\beta^{(N)}\) for (KL.7)
at level \(N\), and \(e_{\beta+\mathbf e_i}^{(N+1)}\) for
the next layer.  The actual recurrence
\(sp_j=p_{j+1}+b_jp_j-(\kappa_j/\kappa_{j-1})p_{j-1}\)
gives the complete retained difference
\[
\begin{split}
 D^{(k)}e_\beta^{(N)}-\sum_i e_{\beta+\mathbf e_i}^{(N+1)}
 ={}&\mathcal T\left[
       \left(\sum_i b_{\beta_i}\right)p_\beta
       -\sum_{i:\beta_i>0}
        \frac{\kappa_{\beta_i}}{\kappa_{\beta_i-1}}
                       p_{\beta-\mathbf e_i}\right]\\
    &-D^{(k)}R_Na_\beta
      +R_{N+1}\sum_i a_{\beta+\mathbf e_i}.
\end{split}
\]
The final term in the one-variable recurrence is omitted at \(j=0\).
Every term on the right has numerator degree at most \(N+1\).
Each term on the left has zero full jet, so their difference is in
\(\mathcal L_{N+1}\).  This proves (KL.20) with all lower
polynomial coefficients retained.  Applying the explicit ordered
division in (KL.9) to this displayed numerator supplies its original
theta primitive.

It follows that \(\partial_N\) is injective and its cokernel is
exactly \(\mathsf H_{N+2}'\), with maps (KL.16)--(KL.17) at
index \(N+1\).  Thus the original derivative connects consecutive
relation quotient types by a proved injection; it has not been made
an endomorphism of a fixed quotient.  The ordinary relation-admission
map has the same declared source and target:
\[
 \operatorname{adm}_N:\mathcal L_{N+1}/\mathcal L_N
       \longrightarrow\mathcal L_{N+2}/\mathcal L_{N+1},
 \qquad [f]\longmapsto[f]=0.
\]
Equation (KL.19) instead sends the representative to its next
derivative class.  These are explicitly different typed maps of the
same original representative, related by (KL.20) and the inclusion
of their actual numerator relations.

All maps here have their original supported lifts.  At a fixed
support label \(\lambda\), a linear map \(T\) acts by
\((\lambda,x)\mapsto(\lambda,Tx)\).  The relation-admission
map therefore sends a killed class to the receiving fibre's
supported zero.  The external zero \(\tau\) remains external
absence.  Each displayed kernel and quotient records its full
supported vector space and each derivative retains its consecutive
source and target; no passage to the Hilbert closure is used.

\subsection{The exact original-theta norm and volume bridge}

Specialize to \(k=1\).  We now distinguish the degree-indexed
\(G_N\) of (KL.6) from the correction-indexed original metric by
writing
\[
 \mathbf G_m=G_{d+m}\quad(m\geq-1),\qquad
 \mathbf R_m=R_{d+m},\qquad
 r_m=\frac{\det\mathbf G_m}{\det\mathbf G_{m-1}}
       \quad(m\geq0).
                                                               \tag{KL.21}
\]
At \(m=-1\), \(\mathcal P_{d-1}\) contains no nonzero
multiple of \(h\).  Thus the original representative is exactly
the fixed section
\[
 s_hu=\mathbf R_{-1}u
   =\mathcal T_h\operatorname{rem}_h(\varepsilon_hu),
 \qquad \varepsilon_h=\upsilon_h^{-1},
 \qquad \mathbf G_{-1}=s_h^*s_h.
\]
This fixes the initial determinant in (KL.21), without introducing
a negative-degree theta polynomial.  Put \(L_{-1}=0\) and
\[
 f_j=D^jf_0,\qquad L_j=\operatorname{span}(f_0,\ldots,f_j),
 \qquad u_j=(I-P_{L_{j-1}})f_j,
 \qquad \mathfrak h_j=\|u_j\|^2>0\quad(j\geq0).
                                                               \tag{KL.22}
\]
These \(u_j=Q_j(D)f_0\) are the monic orthogonal theta
vectors for the original measure
\(|g(1/2+it)|^2dt/(2\pi)\).  Their definition does not depend
on \(h\).  The monic \(p_j\) and their norms \(\kappa_j\)
continue to use the original packet-dependent measure
\(|g/h|^2dt/(2\pi)\).  Multiplication by the unchanged monic
polynomial \(h\) is the exact isometry connecting their source
polynomial spaces:
\[
 P\longmapsto hP,\qquad
 \|hP\|_{\nu_h}^2=\|P\|_{|g|^2dt/(2\pi)}^2,
 \qquad \mathcal T_h(hP)=P(D)f_0.
                                                               \tag{KL.23}
\]
This follows by cancellation of the pointwise factor \(|h|^2\)
in the integral; the identities extend across the discrete zeros by
their entire expressions.  It is an isometry onto the specified
subspace of multiples of \(h\), with its degree shifted by the
full degree \(d=\sum_\rho m_\rho\).

\begin{theorem}[Original norm as an exact relative Gram volume]
For every \(j\geq0\), including the initial theta vector,
\[
 \boxed{\kappa_{d+j}=\mathfrak h_j r_j.}                \tag{KL.24}
\]
Consequently, for \(m\geq0\),
\[
 \boxed{\frac{\kappa_{d+m+1}}{\kappa_{d+m}}
    =\frac{\mathfrak h_{m+1}}{\mathfrak h_m}
                         \frac{r_{m+1}}{r_m},\qquad
 \frac{\det\mathbf G_m}{\det\mathbf G_{-1}}
    =\prod_{j=0}^m\frac{\kappa_{d+j}}{\mathfrak h_j}.} \tag{KL.25}
\]
All ratios have positive denominators and all packet multiplicities
are retained.
\end{theorem}
\begin{proof}
At the degree level \(N=d+j-1\), define
\(e_{d+j}=\mathcal T_hp_{d+j}-\mathbf R_{j-1}a_{d+j}\).
Its numerator has zero full jet, is monic of degree \(d+j\),
and is divisible by the original monic \(h\).  By ordinary monic
division its quotient is monic of degree \(j\).  Under (KL.23)
this is \(f_j\) plus a vector in \(L_{j-1}\).  By (KL.7)
it is orthogonal to that space.  The defining orthogonal projection
in (KL.22) therefore proves
\[
 e_{d+j}=u_j,
 \qquad
 \mathfrak h_j=\kappa_{d+j}
                       +a_{d+j}^*\mathbf G_{j-1}a_{d+j}.
                                                               \tag{KL.26}
\]
When \(j=0\), the quotient of the numerator by \(h\) is the
constant \(1\), so the same argument explicitly gives
\(e_d=f_0=u_0\); no previous theta layer is required.

For the full unit cancellation put
\(v_l=[p_l]_h=p_l(A_h)[1]_h\), and
\(K_{l-1}=\sum_{i=0}^{l-1}v_iv_i^*/\kappa_i\), using the
original monomial remainder coordinates.  At \(l=d+j\),
\[
 a_l=U_{\upsilon_h}v_l,\qquad
 \mathbf G_{j-1}=U_{\varepsilon_h}^*K_{l-1}^{-1}
                                      U_{\varepsilon_h},
 \qquad
 a_l^*\mathbf G_{j-1}a_l=v_l^*K_{l-1}^{-1}v_l.          \tag{KL.27}
\]
Both cancellations use
\(U_{\varepsilon_h}U_{\upsilon_h}=I\); hence they preserve
every derivative of the original local unit at every full jet.
The degree update
\(K_l=K_{l-1}+v_lv_l^*/\kappa_l\) has determinant ratio
\(1+v_l^*K_{l-1}^{-1}v_l/\kappa_l\).  To verify this scalar
identity directly, expand \(\det(I+xy^*)\) multilinearly in
its columns: all terms with two updated columns vanish since those
columns are multiples of \(x\), leaving \(1+y^*x\).
The fixed-unit congruences in (KL.27) then give
\[
 r_j=\frac{\kappa_l}
          {\kappa_l+v_l^*K_{l-1}^{-1}v_l}
     =\frac{\kappa_{d+j}}{\mathfrak h_j}.
\]
This also proves \(0<r_j\leq1\), including the case \(v_l=0\).
Divide two consecutive identities and multiply the first \(m+1\)
identities to obtain (KL.25).
\end{proof}

For two full packets \(h,H\), their original norm sequences are
therefore connected at the same theta degree \(j\) by the exact
scalar equality
\[
 \frac{\kappa_{\deg h+j}^{(h)}}{r_j^{(h)}}
    =\mathfrak h_j
    =\frac{\kappa_{\deg H+j}^{(H)}}{r_j^{(H)}}.           \tag{KL.28}
\]
When \(h\mid H\), the corresponding source isometry is
\(P\mapsto(H/h)P\), since
\((g/H)(H/h)P=(g/h)P\).  At the old centres this equality
preserves every jet, and at each added centre \(g/h\) has the
complete zero order, so it realizes exactly the CRT insertion of
zero jets.  The reference numerator has degree below \(\deg H\)
and is its unique full remainder.  Both corrected representatives
then subtract the same projection onto the packet-independent
\(L_j\), proving the coherent equality
\(\mathbf R_{H,j}\iota_{hH}=\mathbf R_{h,j}\), including
\(j=-1\).  Thus (KL.28) connects precisely these original
coherent source maps and their fixed-section metrics.

\subsection{The complete scalar dictionary for the boundary energy}

For \(m\geq0\), set \(n=d+m+1\) and retain the terminal
notation
\[
 M_n=\kappa_{n-1}^{-1}K_{n-1}^{-1},\quad
 \alpha_n=v_{n-1}^*M_nv_{n-1},\quad
 \beta_n=v_n^*M_nv_n,\quad
 \gamma_n=v_n^*M_nv_{n-1}.
\]
Define the original theta projection rows by
\[
 d_j=\mathfrak h_j^{-1}u_j^*s_h:E_h\to\mathbb C,
 \qquad
 \mathbf R_m=s_h-\sum_{j=0}^m u_jd_j,
 \qquad c_m=d_m\mathbf G_m^{-1}d_{m+1}^*.
                                                               \tag{KL.29}
\]
Here an expression \(u_j^*F\) denotes the Hilbert inner product
\(\langle u_j,F\rangle\), conjugate-linear in its first
argument.  The row \(d_j\) therefore retains the original
unnormalized vector and the exact factor \(\mathfrak h_j^{-1}\).

The complete row and scalar comparison is
\[
\boxed{\begin{aligned}
 d_m&=-\frac{v_{n-1}^*K_{n-1}^{-1}U_{\varepsilon_h}}
                    {\kappa_{n-1}},\\
 d_{m+1}&=-\frac{v_n^*K_{n-1}^{-1}U_{\varepsilon_h}}
                         {\mathfrak h_{m+1}},\\
 d_m\mathbf G_m^{-1}d_m^*
      &=\frac{\alpha_n}{\kappa_{n-1}},\\
 d_{m+1}\mathbf G_m^{-1}d_{m+1}^*
      &=\frac{\kappa_{n-1}\beta_n}{\mathfrak h_{m+1}^2},\\
 \mathfrak h_{m+1}c_m&=\overline{\gamma_n}.
\end{aligned}}                                                 \tag{KL.30}
\]
To prove the first row, the polynomial numerator of \(s_h\)
has degree below \(d\).  In (KL.29), each \(u_j\) has numerator
\(hQ_j\), monic of degree \(d+j\).  Thus the coefficient of
degree \(d+m=n-1\) in the numerator of \(\mathbf R_m\) is
\(-d_m\).  The finite expansion (KL.6) gives the same coefficient
as \(a_{n-1}^*\mathbf G_m/\kappa_{n-1}\).  Cancelling the
same units as in (KL.27) proves the first row, also at \(m=0\).
For the second row use \(u_{m+1}=e_n\), so (KL.10) gives
\(u_{m+1}^*\mathbf R_m=-a_n^*\mathbf G_m\).
The orthogonality of \(u_{m+1}\) to \(L_m\) also gives
\(u_{m+1}^*\mathbf R_m=u_{m+1}^*s_h
 =\mathfrak h_{m+1}d_{m+1}\).  This proves its formula.
Finally
\(\mathbf G_m^{-1}=U_{\upsilon_h}K_{n-1}U_{\upsilon_h}^*\);
insert both rows, cancel both unit pairs and the adjacent kernel
inverse.  The three remaining products are exactly those in
(KL.30), including the conjugation in the last line.

\begin{theorem}[Exact original-norm formula for the reflected packet]
Suppose the original full packet is stable under
\(\rho\mapsto1-\overline\rho\).  For its actual relative
allowance \(\epsilon_m\),
\[
 \boxed{\epsilon_m^2=
  \frac{\mathfrak h_{m+1}}{\mathfrak h_m}
       (r_m^{-1}-1)(1-r_{m+1})
       -\mathfrak h_{m+1}^2|c_m|^2.}                  \tag{KL.31}
\]
If the packet is also stable under ordinary conjugation, then
\(c_m=0\), so its complete formula becomes
\[
 \boxed{\epsilon_m^2=
  \frac{\mathfrak h_{m+1}}{\mathfrak h_m}
                  (r_m^{-1}-1)(1-r_{m+1}).}            \tag{KL.32}
\]
\end{theorem}
\begin{proof}
The terminal identity (KT.4), with the original full multiplicities,
gives
\(\operatorname{Re}\gamma_n=
\operatorname{Re}\operatorname{tr}A_h-d/2=0\): reflected
centres contribute complementary real parts with equal orders.
Its two possible nonzero generalized eigenvalues then yield
\(\epsilon_m^2=\alpha_n\beta_n-|\gamma_n|^2\).
The preceding and succeeding rank-one kernel updates, also proved
directly in (KL.27), give
\[
 \alpha_n=1-r_m,\qquad
 \beta_n=\frac{\kappa_n}{\kappa_{n-1}}
                              \frac{1-r_{m+1}}{r_{m+1}}.
\]
Insert (KL.25) into the product \(\alpha_n\beta_n\) and
the last line of (KL.30) into \(|\gamma_n|^2\).  Every positive
ratio is retained, and this is (KL.31).

For the final assertion, conjugation stability makes \(h\) real
in the original monomial coordinates.  The real Taylor coefficients
of \(g\) give
\((g/h)(1/2-it)=\overline{(g/h)(1/2+it)}\), so the actual
\(\nu_h\) is even in \(t\).  Its monomial Gram matrix is real:
the integrand for entry \(\overline{s^i}s^j\) at \(-t\) is
the complex conjugate of its integrand at \(t\).  Therefore the
unique solution of the real Gram linear equations for the monic
\(p_l\) has real coefficients.  Polynomial division by the real
monic \(h\) makes each \(v_l\) real, and the finite sum for
\(K_{n-1}\) is real.  Its inverse and \(M_n\) are real as well.
Consequently \(\gamma_n\) is real.  The already proved reflection
trace identity makes its real part zero, hence \(\gamma_n=0\).
Equation (KL.30), with \(\mathfrak h_{m+1}>0\), proves
\(c_m=0\).  No division by a jet column, an off-line displacement,
or a boundary energy has been used; all zero-column cases remain
included.
\end{proof}

The source PR15 denotes the relative determinant at correction
\(m\) by \(\delta_m=\det\mathbf G_{m+1}/\det\mathbf G_m\).
Its exact index map is therefore \(\delta_m=r_{m+1}\), with
\(\delta_{-1}=r_0\).  Its cross-term row product is the
\(c_m\) of (KL.29).  Equations (KL.24)--(KL.31) prove the full
map from its original-theta scalar formula to the terminal-kernel
formula, including both monic norm families, the unit congruences,
the conjugation, the fixed initial section, and the complete orders.
These equalities calculate the same actual boundary energy at every
finite stage.  They do not supply an estimate tending to zero; that
requires further information about these now explicitly related
arithmetic quantities.


The packet-independent norms in (KL.24)--(KL.32) also have the
complete positive-half-line formulas (SP.18)--(SP.21) and the
gamma-reference determinant formulas (TG.16)--(TG.19). Those
sections prove the original affine and quadratic coordinate maps,
their inverses, both measures and every phase. Consequently these
formulas can be substituted directly into the identities here;
the actual unit, packet determinant ratios, full zero orders and
source degree indices remain unchanged.
